\documentclass[11pt]{article}
\usepackage[T1]{fontenc}
\usepackage{lmodern,microtype,amsmath,amssymb,amsthm,booktabs,longtable,array}
\usepackage[margin=1in]{geometry}
\usepackage[colorlinks=true,linkcolor=blue,citecolor=blue,urlcolor=blue]{hyperref}
\usepackage{enumitem}
\setlist{nosep}
\newtheorem{theorem}{Theorem}[section]
\newtheorem{lemma}[theorem]{Lemma}
\newtheorem{proposition}[theorem]{Proposition}
\newtheorem{corollary}[theorem]{Corollary}
\newtheorem{remark}[theorem]{Remark}
\newcommand{\Z}{\mathbb Z}
\newcommand{\Q}{\mathbb Q}
\newcommand{\C}{\mathbb C}
\newcommand{\Hh}{\mathbb H}
\DeclareMathOperator{\trdeg}{trdeg}
\DeclareMathOperator{\ord}{ord}
\title{Effective CM Nonvanishing for Products of Eisenstein Series\\[4pt]\large Infinite Weight Bands, a Finite Reduction, and a Two-Adic Obstruction}
\author{Research note prepared for Keio}
\date{14 November 2026}
\begin{document}
\maketitle
\begin{abstract}
Let $F_{k,\ell}=E_kE_\ell-E_{k+\ell}$, where the Eisenstein series are normalized to have constant term one. We give an effective finite reduction for nonvanishing at each fixed principal CM point of odd discriminant $D\le-7$. A separate, reproducible interval computation and a uniform lattice estimate establish nonvanishing at all 35 principal CM points in the boundary classification whenever $\min(k,\ell)\le500$, with the identically zero weights excluded. We also prove a two-adic congruence for every even $k,\ell\ge6$ with $v_2(k)=v_2(\ell)$, excluding the CM discriminants $-7,-15$ throughout that infinite family. Given the no-Siegel-zero input used in the boundary classification, the weight-band result implies the corresponding part of Aiken--Ross--Shvydkoy--Xue Conjecture 5.2. The complete conjecture is not proved: the effective global cutoff is not evaluated or exhaustively checked here. No claim of historical priority or formal proof-assistant verification is made.
\end{abstract}

\section{Scope and external inputs}
For an even integer $n\ge4$, put
\[
 E_n(z)=1+c_n\sum_{m\ge1}\sigma_{n-1}(m)q^m,
 \qquad c_n=-\frac{2n}{B_n},\qquad q=e^{2\pi iz}.
\]
All weight pairs in this note are even and at least four. We exclude $k+\ell\in\{8,10,14\}$, for which $F_{k,\ell}$ is identically zero.

The target conjecture states that the zeros of $F_{k,\ell}$ in the standard fundamental domain, apart from $i$ and $\rho=(-1+i\sqrt3)/2$, are transcendental \cite{ARSX}. We use the boundary-zero and CM-conjugacy results assembled in that paper. Under its no-Siegel-zero hypothesis, the only nonelliptic CM discriminants that can occur are
\begin{align*}
\mathcal D=\{&-7,-11,-19,-27,-43,-67,-163,\\
&-15,-35,-51,-75,-91,-99,-115,-123,-147,-187,-235,-267,-403,-427,\\
&-195,-315,-435,-483,-555,-595,-627,-715,-795,-1435,\\
&-1155,-1995,-3003,-3315\}.
\end{align*}
There are 35 discriminants and 107 nonelliptic CM classes. Rational Fourier coefficients imply that a CM zero forces a zero at every CM point of the same discriminant. It therefore suffices to exclude the 35 principal representatives
\[
 \tau_D=\frac{-1+i\sqrt{|D|}}2\qquad(D\in\mathcal D).
\]
The fixed-point nonvanishing theorems below do not require any hypothesis about Dirichlet $L$-functions. The no-Siegel-zero input is used only to assert that this CM list is complete. In the motivating application it is supplied by the proposed input labelled 003; this note does not independently audit that input.

\begin{theorem}[Certified infinite band]\label{thm:band}
For every $D\in\mathcal D$, every admissible weight pair with $\min(k,\ell)\le500$ satisfies
\[
 F_{k,\ell}(\tau_D)\ne0.
\]
Consequently, under the complete boundary-CM classification above, Conjecture 5.2 holds for all these weight pairs.
\end{theorem}

This is an infinite family, not merely a check of total weight at most 1000: the larger weight is unrestricted. The finite triangle $k+\ell\le1000$ was already reported verified by the original authors, along with a stronger polynomial assertion. Our calculation independently checks the nonvanishing needed here; it does not independently check their irreducibility assertion. The cases with smaller weight $4,6,8$ also have a previously known lower-arc argument \cite{RVY}. They are included for completeness, not presented as new.

\section{Effective reduction to finitely many weights}
\label{sec:finite-weight-reduction}

The result of this section concerns a fixed CM point and is unconditional.
The exclusion of Siegel zeros is needed only afterwards, to replace the
possible CM discriminants in the modular-form problem by a known finite list.
An effectively computable bound below is not a numerical bound already
evaluated: no assertion that it lies within the previously tested weights is
made.

\begin{theorem}[Effective finite-weight reduction]
\label{thm:effective-finite-weight}
Let $D\leq -7$ be an odd quadratic discriminant, and put
\[
 z=\frac{-1+i\sqrt{|D|}}2,\qquad
 y=\Im z,\qquad a=\frac{|D|+1}{4},\qquad
 r=\sqrt a,\qquad R=\sqrt{a+2}.
\]
There is an effectively computable integer $N_D$ such that, for even
integers $k,\ell\geq4$,
\[
 \max(k,\ell)>N_D
 \quad\Longrightarrow\quad
 E_k(z)E_\ell(z)-E_{k+\ell}(z)\ne0.
\]
Here $E_n$ denotes the normalized level-one Eisenstein series with constant
term $1$.  Thus the pairs of weights giving equality at this fixed point
form a finite set contained in a computable box.
\end{theorem}

\subsection{The primitive lattice expansion and its error}

The normalization is important.  For even $n\geq4$ one has
\begin{equation}
\label{eq:fr-primitive}
 E_n(z)=\frac12\sum_{\substack{(c,d)\in\mathbb Z^2\\(c,d)=1}}
 (cz+d)^{-n}
 =1+\sum_{\substack{c\geq1,\ d\in\mathbb Z\\(c,d)=1}}(cz+d)^{-n}.
\end{equation}
The two terms $(c,d)=(1,0),(1,1)$ are respectively $z^{-n}$ and
$\bar z^{-n}$, since $z+1=-\bar z$ and $n$ is even.  Their squared
length is $a$.  Every other vector in the sum has squared length at least
$a+2$.  Indeed,
\[
 |cz+d|^2=d^2-cd+ac^2=(d-c/2)^2+c^2(a-1/4).
\]
For $c=1$ the next choices $d=-1,2$ give $a+2$; for $c\geq2$,
the lower bound $c^2(a-1/4)\geq4a-1\geq a+2$ applies.

We record an explicit summable majorant.  By comparison with the integral
of the even decreasing function $(t^2+(cy)^2)^{-2}$ on either side of its
maximum,
\[
 \sum_{d\in\mathbb Z}|cz+d|^{-4}
 \leq 2(cy)^{-4}+\int_{\mathbb R}(t^2+(cy)^2)^{-2}\,dt
 =2(cy)^{-4}+\frac\pi2(cy)^{-3}.
\]
Consequently, because $y\geq\sqrt7/2>13/10$,
\begin{equation}
\label{eq:fr-fourth-moment}
 \sum_{c\geq1}\sum_{d\in\mathbb Z}|cz+d|^{-4}
 \leq 2\zeta(4)y^{-4}+\frac\pi2\zeta(3)y^{-3}
 <\frac4{y^3}.
\end{equation}
For example, the last inequality follows already from
$\zeta(4)<11/10$, $\zeta(3)<5/4$, and $\pi<22/7$.

Define
\[
 u_n=z^{-n}+\bar z^{-n},\quad A_n=E_n(z)-1,\quad
 B_D=\frac{4a^2}{y^3},\quad Q_D=\frac{4(a+2)^2}{y^3}.
\]
The norm gap and \eqref{eq:fr-fourth-moment} imply, uniformly for even
$n\geq4$,
\begin{equation}
\label{eq:fr-expansion}
 |A_n|\leq B_Dr^{-n},\qquad
 A_n=u_n+\varepsilon_n,\qquad
 |\varepsilon_n|\leq Q_DR^{-n}.
\end{equation}
This is a primitive-lattice expansion, so no factor $\zeta(n)$, or error
arising from such a factor, is being suppressed.

By symmetry we may assume $k\geq\ell$.  Put
$F_{k,\ell}(z)=E_k(z)E_\ell(z)-E_{k+\ell}(z)$.  From
\eqref{eq:fr-expansion},
\begin{equation}
\label{eq:fr-main-error}
 F_{k,\ell}(z)=u_k+u_\ell+\mathcal R_{k,\ell},\qquad
 |\mathcal R_{k,\ell}|
 \leq2Q_DR^{-\ell}+(B_D^2+B_D)a^{-\ell}
 \leq C_Db^{-\ell},
\end{equation}
where $b=\min(a,R)>r$ and one may take
$C_D=2Q_D+B_D^2+B_D$.

\subsection{An arithmetic obstruction to exact cancellation}

Let $S_n=z^n+\bar z^n$.  Since $z^2+z+a=0$,
\[
 S_0=2,\qquad S_1=-1,\qquad
 S_n=-S_{n-1}-aS_{n-2}.
\]
In particular, $S_n\in\mathbb Z$ and
\begin{equation}
\label{eq:fr-trace-congruence}
 S_n\equiv(-1)^n\pmod a\qquad(n\geq1).
\end{equation}
As $u_n=S_n/a^n$, if $k>\ell$ are even then
\[
 a^k(u_k+u_\ell)=S_k+a^{k-\ell}S_\ell\equiv1\pmod a.
\]
If $k=\ell$, the same sum is $2S_\ell/a^\ell\ne0$, since
$S_\ell\equiv1\pmod a$.  We have therefore proved
\begin{equation}
\label{eq:fr-main-nonzero}
 u_k+u_\ell\ne0\qquad(k,\ell\geq4\text{ even}).
\end{equation}
The quantitative argument next requires precisely this exclusion of exact
cancellation.

\subsection{A uniform lower bound for the two leading terms}

Set $\beta=z/\bar z$.  This algebraic number has modulus $1$ and is not
a root of unity.  Indeed,
\[
 \beta+\beta^{-1}=\frac1a-2\in\mathbb Q\setminus\mathbb Z.
\]
If $\beta$ were a root of unity, this rational number would be an
algebraic integer, which is impossible.

We use the effective lower bound for a nonzero linear form in logarithms
of algebraic numbers of Baker--W\"ustholz \cite{BW}.  In the present application
there are at most three logarithms, all their arguments lie in the fixed
quadratic field $\mathbb Q(z)$, and all have modulus $1$ in the chosen
complex embedding.  Thus the usual modified heights can be bounded by
constant multiples of $\max(1,h(\alpha))$, where $h$ is the absolute
logarithmic height.  The resulting constants are effective.

First apply this theorem to the logarithms of $\beta$ and $-1$.
There are effective constants $0<c_D'\leq1$ and $A_D'>0$ such that
\begin{equation}
\label{eq:fr-one-trace}
 |u_\ell|=r^{-\ell}|1+\beta^\ell|
 \geq c_D'r^{-\ell}(\ell+2)^{-A_D'}
 \qquad(\ell\geq4\text{ even}).
\end{equation}
The linear form cannot vanish, because $\beta$ is not a root of unity.

Write $m=k-\ell\geq0$.  If
\begin{equation}
\label{eq:fr-large-gap}
 m\geq\frac{A_D'\log(\ell+2)+\log(4/c_D')}{\log r},
\end{equation}
then $|u_k|\leq2r^{-k}\leq\tfrac12|u_\ell|$.  Hence
\[
 |u_k+u_\ell|
 \geq\frac{c_D'}2r^{-\ell}(\ell+2)^{-A_D'}.
\]

It remains to consider gaps not satisfying \eqref{eq:fr-large-gap}.
There is an effective constant $H_D$ such that all these gaps satisfy
$m\leq H_D\log(\ell+2)$.  Define
\[
 \delta_m=\frac{1+\bar z^{-m}}{1+z^{-m}}.
\]
Both its numerator and denominator are nonzero.  For $m=0$ the two
are $2$; for positive even $m$ one has
$|1+z^{-m}|\geq1-r^{-2}\geq1/2$.  We have
\begin{equation}
\label{eq:fr-two-trace-factorization}
 u_k+u_\ell
 =z^{-\ell}(1+z^{-m})\bigl(1+\beta^\ell\delta_m\bigr).
\end{equation}
Here $\delta_m\in\mathbb Q(z)$ and $|\delta_m|=1$.
Since $h(z)=\log r$, the standard height inequalities give the explicit
estimate
\begin{equation}
\label{eq:fr-height}
 h(\delta_m)\leq2m\log r+2\log2
 =m\log a+2\log2.
\end{equation}

For completeness, take principal logarithms of the unit-modulus numbers
$\beta$ and $\delta_m$.  Choose an odd integer $v$ so that
\[
 \Lambda=\ell\log\beta+\log\delta_m-v\log(-1)
\]
is purely imaginary with imaginary part in $[-\pi,\pi]$.
Then $|v|\leq\ell+2$, and
$|\exp\Lambda-1|=|1+\beta^\ell\delta_m|$.
By \eqref{eq:fr-main-nonzero}, $\Lambda\ne0$.  An argument equal to
$1$ can simply be omitted from the linear form.  Applying the
three-logarithm bound, using \eqref{eq:fr-height}, and then using
$|e^{it}-1|\geq2|t|/\pi$ for $|t|\leq\pi$, yields
\[
 |1+\beta^\ell\delta_m|
 \geq \exp\{-K_D(m+1)\log(\ell+3)\}
\]
after increasing an effective positive constant $K_D$ if necessary.
Together with $m\leq H_D\log(\ell+2)$ and
\eqref{eq:fr-two-trace-factorization}, this gives an effective lower
bound of the form
\[
 |u_k+u_\ell|
 \geq c_D''r^{-\ell}
 \exp\{-A_D''(\log(\ell+2))^2\}.
\]
The bound from the large-gap case is at least as strong after enlarging
$A_D''$ and decreasing $c_D''$.  Consequently this last inequality
holds \emph{uniformly for every even $k\geq\ell\geq4$}.

Since $b/r>1$, there is an effectively computable $L_D$ such that
\[
 c_D''r^{-\ell}e^{-A_D''(\log(\ell+2))^2}
 >C_Db^{-\ell}\qquad(\ell\geq L_D).
\]
For example, choose $L_D$ beyond a point where
\[
 \ell\log(b/r)-A_D''(\log(\ell+2))^2
 +\log(c_D''/C_D)
\]
is positive and increasing.  Such a point can be found effectively,
because its linear term eventually dominates its squared logarithm.
Equation \eqref{eq:fr-main-error} now proves nonvanishing whenever
$\min(k,\ell)\geq L_D$.

\subsection{Removing the remaining infinite columns}

We still must handle the finitely many even $\ell$ with
$4\leq\ell<L_D$ and arbitrarily large $k$.  The relevant limiting
value is $E_\ell(z)-1$, and it is essential to prove that this value
does not vanish.

The point $z$ has algebraic $j$-invariant because it is CM.  Moreover,
$E_4(z)E_6(z)\ne0$: in the standard fundamental domain the zeros of
$E_4$ and $E_6$ occur only at the elliptic points, whereas $\Im z>1$.
Put
\[
 w=\frac{E_6(z)}{E_4(z)},\qquad
 c=\frac{E_4(z)^3}{E_6(z)^2}
   =\frac{j(z)}{j(z)-1728}\in\overline{\mathbb Q}^{\times}.
\]
Thus $E_4(z)=cw^2$ and $E_6(z)=cw^3$.  The graded-ring identity
$M_*({\rm SL}_2(\mathbb Z);\mathbb Q)=\mathbb Q[E_4,E_6]$
implies
\[
 E_\ell(z)=\lambda_\ell w^{\ell/2},\qquad
 \lambda_\ell\in\overline{\mathbb Q}.
\]
If $E_\ell(z)=1$, then $\lambda_\ell\ne0$ and $w$ would be
algebraic.  Hence $E_4(z)$ and $E_6(z)$ would both be algebraic.
This contradicts Nesterenko's theorem \cite{Nesterenko}
\[
 \operatorname{trdeg}_{\mathbb Q}
 \mathbb Q(q,E_2(z),E_4(z),E_6(z))\geq3,
 \qquad q=e^{2\pi iz},\quad 0<|q|<1.
\]
We conclude that $E_\ell(z)\ne1$ for every even $\ell\geq4$.

This qualitative nonvanishing suffices for effectivity.  For each of
the finitely many $\ell<L_D$, the absolutely convergent Fourier series,
with an explicit tail bound and certified approximations to $q$, gives
successively smaller rational intervals for the real number
$E_\ell(z)-1$.  Because this number is nonzero, eventually such an
interval excludes $0$.  We thereby find a positive rational number
$\eta_\ell<|E_\ell(z)-1|$ and a rational upper bound
$M_\ell>|E_\ell(z)|$.

Finally, \eqref{eq:fr-expansion} gives
\[
 \left|F_{k,\ell}(z)-(E_\ell(z)-1)\right|
 \leq B_D(M_\ell+1)r^{-k}.
\]
Choose an effective $K_\ell$ such that the right-hand side is less
than $\eta_\ell$ whenever $k\geq K_\ell$.  Taking $N_D$ to be
an integer at least $L_D$ and at least every such $K_\ell$ proves
Theorem~\ref{thm:effective-finite-weight}.  This is a terminating
algorithm for a bound, even though its running time and its output may
be very large.

\subsection{Consequence for the global transcendence question}

Assume the absence of Siegel zeros, and use the established boundary
location and CM-exception results for $F_{k,\ell}$.  The nonelliptic
CM exceptions then have one of the $35$ odd discriminants
$D\ne-3$ whose order class group has exponent dividing $2$.
If $F_{k,\ell}$ vanishes at any CM point of discriminant $D$, its
rational Fourier coefficients imply that it vanishes at the principal
representative $\tau_D$ as well.  Thus, taking
\[
 N=\max_{D\in\mathcal D}N_D,
\]
Theorem~\ref{thm:effective-finite-weight} shows that every nonelliptic
zero is transcendental whenever $\max(k,\ell)>N$.
The only remaining weight pairs lie in the computable finite box
$4\leq k,\ell\leq N$ (with even weights and
$k+\ell\notin\{8,10,14\}$).

\paragraph{Scope of the conclusion.}
The numerical value of $N$ has not been evaluated here, and the resulting
finite box has not been exhaustively checked.  In particular, no comparison
between $N$ and the previously reported verification range
$k+\ell\leq1000$ is established.  This section proves an effective
finite-weight reduction, not the complete transcendence conjecture.


\section{An explicit infinite weight band}
\subsection{A uniform estimate beyond the finite triangle}
Write $R_n=E_n(\tau_D)-1$, $y=\sqrt{|D|}/2$ and $a=(|D|+1)/4$. The primitive lattice estimate proved above gives
\begin{equation}\label{eq:bandbound}
 |R_n|<4y^{-3}a^{2-n/2}\qquad(n\ge4).
\end{equation}
For $n\ge10$ and $a\ge2$, this also gives $|R_n|<1$ and hence $|E_n|<2$. If $\ell\ge10$ and $k\ge K$, then
\begin{align}
 |F_{k,\ell}-(E_\ell-1)|
 &=|E_\ell R_k-R_{k+\ell}|\notag\\
 &<12y^{-3}a^{2-k/2}
 \le12y^{-3}a^{2-K/2}.\label{eq:domination}
\end{align}
For every even $10\le\ell\le496$, set $K=1000-\ell$. The certified computation verifies, separately at all 35 discriminants,
\begin{equation}\label{eq:finitebandtest}
 |R_\ell|>12y^{-3}a^{2-K/2}.
\end{equation}
There are $244\cdot35=8540$ successful tests. The smallest ratio of the left-hand side to the right-hand side occurs at $(\ell,D)=(496,-7)$ and exceeds $1.5218$. Thus \eqref{eq:domination} proves nonvanishing for every $k\ge1000-\ell$. The remaining pairs have total weight at most 1000 and are covered by the independent finite-triangle verification below. This proves the band through smaller weight 496.

For the two final columns, use the same inequality with $K_{498}=508$ and $K_{500}=516$. All 70 tests pass, with ratios exceeding $1.9427$ and $1.0364$, respectively. The finite gaps outside the original triangle are exactly the nine unordered pairs
\[
 (498,504),\ (498,506),\quad (500,k)\quad(k=502,504,\ldots,514).
\]
Their 315 evaluations at the 35 points also exclude zero. Together with the known smaller weights $4,6,8$, this proves Theorem~\ref{thm:band}. The initial cutoff $K=1000-\ell$ was insufficient for these last two columns; increasing $K$ and checking the finite gaps resolves that limitation.

\subsection{Certified evaluation and finite enumeration}
The principal points have real negative $q=-r$, where $r=e^{-\pi\sqrt{|D|}}<1/4000$. For $s=w-1\ge3$, write
\[
 S_{s,N}(q)=\sum_{n=1}^{N}\sigma_s(n)q^n.
\]
The coefficients are exact integers and the Bernoulli normalization is an exact rational number. Since $\sigma_s(n)\le\zeta(s)n^s<2n^s$, the omitted tail is bounded by
\begin{equation}\label{eq:tail}
 \left|\sum_{n>N}\sigma_s(n)q^n\right|
 \le\frac{2(N+1)^s r^{N+1}}{1-r((N+2)/(N+1))^s},
\end{equation}
provided the denominator is positive. This holds for $N=300$ and every $4\le w\le1014$ in use here; the scripts also check it by interval arithmetic.

We evaluate the truncated polynomial by Horner's rule, add the symmetric interval enclosing \eqref{eq:tail}, and multiply by the exact rational $c_w$. All certification decisions use outward interval arithmetic at 400 decimal digits in \texttt{mpmath.iv}. The auxiliary ordinary multiprecision context is used only for printing and ranking, never for a nonvanishing decision. For each discriminant, cached intervals $R_w$ yield
\[
 F_{k,\ell}=R_k+R_\ell+R_kR_\ell-R_{k+\ell}.
\]
This avoids subtracting two quantities with constant term one.

The finite domain is precisely
\[
 4\le k\le\ell,\quad k,\ell\text{ even},\quad k+\ell\le1000,
 \quad k+\ell\notin\{8,10,14\}.
\]
It contains 61,997 unordered weight pairs and therefore 2,169,895 point--pair evaluations over the 35 discriminants. The executable verifier requires each resulting interval to be strictly positive or strictly negative.

\begin{center}
\begin{tabular}{@{}lll@{}}
\toprule
Verification & Domain & Purpose\\
\midrule
Finite triangle & 2,169,895 evaluations & All total weights $\le1000$\\
Infinite-band tests & 8,610 inequalities & All larger weights, $10\le\min\le500$\\
Edge gaps & 315 evaluations & Nine additional weight pairs\\
Small-column check & 2,065 evaluations & Separate check for weights 10 and 12\\
\bottomrule
\end{tabular}
\end{center}
The package contains the scripts, JSON summaries and the band-test records. Printed decimal endpoints in diagnostic JSON fields are rounded summaries; certification is performed with the internal outward intervals. Re-running the scripts reconstructs those intervals from exact integer inputs and the explicit analytic tail. This is a reproducible computational proof component with a numerical-library trust base, not a Lean certificate.

% Self-contained mathematical appendix.  This file is intended to be
% \input into a document that loads amsmath, amssymb, and amsthm.
\section{A two-adic nonvanishing theorem for equal valuations of the weights}
\label{sec:two-adic-equal-valuations}

Write $\mathbb Z_{(2)}$ for the rational numbers with odd denominator and
normalize the valuation by $v_2(2)=1$.  Congruences of rational numbers
are understood in $\mathbb Q_2$; congruences of Fourier expansions are
coefficientwise.  We use the normalization
\[
 E_n(q)=1+a_nS_n(q),\qquad
 a_n=-\frac{2n}{B_n},\qquad
 S_n(q)=\sum_{m\ge1}\sigma_{n-1}(m)q^m
\]
for every even $n\ge4$.  As usual,
$\Delta=q\prod_{m\ge1}(1-q^m)^{24}$ and $j=E_4^3/\Delta$.

\paragraph{Theorem.}
Let $k,\ell\ge6$ be even integers with
$v_2(k)=v_2(\ell)$.  Set $w=k+\ell$ and
\[
 G_{k,\ell}=\frac{E_kE_\ell-E_w}{a_ka_\ell}.
\]
Then $G_{k,\ell}$ has coefficients in $\mathbb Z_{(2)}$, and
\begin{equation}
 G_{k,\ell}\equiv
 \Delta E_4^{(w-12)/4}\pmod2.
 \label{eq:equal-v-two-congruence}
\end{equation}
In particular, $w$ is divisible by $4$.  If
$N=\lfloor w/12\rfloor$ and $b=w/4-3N\in\{0,1,2\}$, then
\[
 G_{k,\ell}=\Delta^NE_4^bP_{k,\ell}(j)
\]
for a polynomial $P_{k,\ell}\in\mathbb Z_{(2)}[X]$ satisfying
\begin{equation}
 P_{k,\ell}(X)\equiv X^{N-1}\pmod2.
 \label{eq:equal-v-two-j-polynomial}
\end{equation}
Consequently, a point $\tau$ whose algebraic $j$-invariant is a
$2$-adic unit at at least one prime above $2$ cannot be a zero of
$E_kE_\ell-E_w$.

\paragraph{Bernoulli lemma.}
For every even $n\ge6$, if $r=v_2(n)$, then
\begin{equation}
 2B_n\equiv1+2n\pmod{2^{r+3}},\qquad
 a_n\equiv-4n+8n^2\pmod{2^{2r+5}}.
 \label{eq:bernoulli-two-precision}
\end{equation}
Here is an elementary proof of the precision needed below.  The
duplication identity for Bernoulli polynomials gives
\[
 2(1-2^n)B_n
 =1-n+\sum_{\substack{2\le h\le n-2\\h\text{ even}}}
       \binom nh2^hB_h.
\]
The von Staudt--Clausen theorem gives $v_2(B_h)=-1$ for positive even
$h$, while
\[
 v_2\binom nh\ge r-v_2(h)
\]
follows from $\binom nh=(n/h)\binom{n-1}{h-1}$.  For every even
$h\ge6$ the corresponding summand therefore has valuation at least
$r+h-1-v_2(h)\ge r+4$.  Retaining $h=2,4$ yields
\[
 2B_n\equiv
 1-n+\frac{n(n-1)}3
 -\frac{n(n-1)(n-2)(n-3)}{45}
 \pmod{2^{r+3}}.
\]
We used $n\ge r+3$ to discard the factor $2^n(2B_n)$.  The expression
on the right equals
\[
 1-\frac{6n}{5}+\frac{4n^2}{45}
       +\frac{2n^3}{15}-\frac{n^4}{45}.
\]
Its difference from $1+2n$ is
\[
 -\frac{16n}{5}+\frac{4n^2}{45}
       +\frac{2n^3}{15}-\frac{n^4}{45}.
\]
The four terms have valuations at least $r+4$, $2r+2$, $3r+1$, and
$4r$, respectively, all at least $r+3$ because $r\ge1$.  This proves
the first congruence in \eqref{eq:bernoulli-two-precision}.  Since
$a_n=-4n/(2B_n)$, inversion gives the second congruence: the error
from the first congruence has valuation at least $2r+5$, and the
first omitted geometric-series term has valuation $3r+4\ge2r+5$.

\paragraph{Proof of the Fourier congruence.}
Put $r=v_2(k)=v_2(\ell)$, $t=v_2(w)$,
$a=a_k$, $b=a_\ell$, and $c=a_w$.  Then $t\ge r+1$, and
von Staudt--Clausen gives
\[
 v_2(a)=v_2(b)=r+2,\qquad v_2(c)=t+2.
\]
Applying the Bernoulli lemma to the three weights gives
\[
 a+b-c\equiv
 8(k^2+\ell^2-w^2)=-16k\ell
 \pmod{2^{2r+5}}.
\]
Thus
\begin{equation}
 u:=\frac{a+b-c}{ab}\in\mathbb Z_{(2)}^\times,
 \qquad u\equiv1\pmod2.
 \label{eq:first-coeff-unit}
\end{equation}
The Fourier expansion is
\[
 G_{k,\ell}=S_kS_\ell+L,\qquad
 [q^n]L=
 \frac{a\sigma_{k-1}(n)+b\sigma_{\ell-1}(n)
                    -c\sigma_{w-1}(n)}{ab}.
\]

We first take $n$ odd.  For each odd divisor $d$ of $n$,
subtracting $(a+b-c)d^{w-1}$ from its contribution to the numerator
leaves
\[
 a d^{k-1}(1-d^\ell)+b d^{\ell-1}(1-d^k).
\]
For any odd $d$, $d^{2^r}\equiv1\pmod{2^{r+2}}$; this is the
usual elementary lifting-the-exponent congruence for odd integers.
Furthermore, because $k/2^r$ and $\ell/2^r$ are both odd, the
geometric-sum identity gives
\[
 \frac{d^k-1}{2^{r+2}}
 \equiv\frac{d^\ell-1}{2^{r+2}}\pmod2.
\]
These quotients are integral, including for $d=1$ where both are zero.
Since $a/2^{r+2}$ and $b/2^{r+2}$ are odd, the two displayed
correction terms cancel modulo $2^{2r+5}$.  Dividing by $ab$ and
using \eqref{eq:first-coeff-unit}, we find
\begin{equation}
 [q^n]L\equiv\sigma_{w-1}(n)\pmod2
 \qquad(n\text{ odd}).
 \label{eq:linear-odd-coefficients}
\end{equation}

For general $n=2^vm$ with $m$ odd, one has
\[
 \sigma_{j-1}(2^vm)=\sigma_{j-1}(m)
       (1+2^{j-1}+\cdots+2^{v(j-1)}).
\]
The resulting corrections after division by $ab$ have valuations
at least
\[
 k-r-3,\qquad \ell-r-3,\qquad w+t-2r-3,
\]
respectively.  All are at least $1$: for the first two use
$k,\ell\ge6$ and $r=v_2(k)=v_2(\ell)$; for the third use
$t\ge r+1$ and $w\ge k+6$.  Hence
\begin{equation}
 [q^n]L\equiv\sigma_{w-1}(m)\pmod2.
 \label{eq:linear-all-coefficients}
\end{equation}
In particular every coefficient of $L$ is $2$-integral.

Modulo $2$, both $S_k$ and $S_\ell$ equal the series $S$ whose
coefficient at $n$ is $1$ exactly when the odd part of $n$ is a square.
This is the parity formula for the number of odd divisors of $n$.
In characteristic $2$ their product is $S^2=S(q^2)$.  Its
coefficients cancel those of $L$ at every even index in
\eqref{eq:linear-all-coefficients}; at odd indices only odd squares
remain.  Therefore
\[
 G_{k,\ell}\equiv\sum_{h\ge0}q^{(2h+1)^2}\equiv\Delta\pmod2.
\]
For completeness, the last identity follows from Jacobi's identity
\[
 \prod_{m\ge1}(1-q^m)^3
 =\sum_{h\ge0}(-1)^h(2h+1)q^{h(h+1)/2}.
\]
In characteristic $2$, substitute $q^8$ on the right and note that
$\Delta\equiv q\prod_{m\ge1}(1-q^{8m})^3$.
Because $E_4\equiv1\pmod2$, the congruence with equal weights is
exactly \eqref{eq:equal-v-two-congruence}.

\paragraph{Proof of the polynomial assertion and nonvanishing.}
Write $w=4m$, $N=\lfloor m/3\rfloor$, and $b=m-3N$.
A triangular basis for weight $4m$ modular forms is
\[
 E_4^{m-3h}\Delta^h,\qquad 0\le h\le N.
\]
Each basis vector begins with $q^h$ with coefficient $1$ and has
integer Fourier coefficients.  Successively comparing these initial
coefficients shows that a form with coefficients in $\mathbb Z_{(2)}$
has basis coefficients in $\mathbb Z_{(2)}$.  For $G_{k,\ell}$ the
coefficient at $h=0$ is zero, and
\eqref{eq:equal-v-two-congruence} says that the coefficient at $h=1$
is odd while those at $h\ge2$ are even.  Factoring out
$\Delta^NE_4^b$ now gives
\eqref{eq:equal-v-two-j-polynomial}.  If $j(\tau)$ is a unit at a
prime above $2$, then its substitution into the reduction of $P$
is nonzero, so $P(j(\tau))\ne0$.  Also $E_4(\tau)\ne0$, since
$j(\tau)\ne0$, and $\Delta$ has no zero in the upper half-plane.
This proves the theorem.

\paragraph{Application to the two ordinary discriminants in the CM list.}
The Hilbert class polynomials
\[
 H_{-7}(X)=X+3375,\qquad
 H_{-15}(X)=X^2+191025X-121287375
\]
are monic with odd constant terms, so all of their roots are units
at every prime above $2$.  Consequently, at every CM point of
discriminant $-7$ or $-15$,
\[
 E_k(\tau)E_\ell(\tau)-E_{k+\ell}(\tau)\ne0
\]
for all even $k,\ell\ge6$ with $v_2(k)=v_2(\ell)$.  This includes
the entire diagonal $k=\ell\ge6$.

\paragraph{Scope.}
This is a uniform theorem for an infinite two-parameter set of
weights, not a finite computational verification.  The proof does
not use the quasi-Riemann hypothesis.  Its role in the larger
argument is to eliminate the $D=-7,-15$ cases from the finite CM
exception list obtained using the zero-free input.  It does not
eliminate the other $33$ discriminants, does not cover weights of
unequal $2$-adic valuation, and does not prove the complete
transcendence conjecture.  No priority claim is made for this
congruence or its consequence.


\section{What remains for the full conjecture}
Let $N_D$ be the effective cutoff from the finite-reduction theorem, and put $N=\max_{D\in\mathcal D}N_D$. Under the complete boundary-CM classification, any still-possible counterexample must have
\[
 502\le\min(k,\ell)\le\max(k,\ell)\le N,
 \qquad k+\ell>1000.
\]
At the discriminants $-7,-15$, the equal-$v_2$ family proved in Section~\ref{sec:two-adic-equal-valuations} is excluded as well. Neither a numerical value of $N$ nor a verification of this remaining region is supplied here. Logarithmic-form constants may make the resulting cutoff very large. The reduction therefore proves effective finiteness, not practical tractability or emptiness of the finite remainder.

Three mathematically distinct ingredients have been combined: the boundary geometry and CM classification restrict the locations; exact quadratic arithmetic prevents cancellation of the shortest-vector main term; logarithmic-form estimates make that noncancellation quantitatively stronger than the exponential lattice tail. The two-adic congruence gives a separate algebraic exclusion. None of these steps assumes the stronger irreducibility conjecture.

\begin{thebibliography}{9}
\bibitem{ARSX} D. Aiken, E. Ross, D. Shvydkoy and H. Xue,
\emph{Transcendence of zeros of modular forms}, author-hosted preprint, submitted September 2025.
\url{https://erick-ross.github.io/transcendence-of-zeros.pdf}.
See Theorem 1.2, Lemma 2.2, Lemma 3.1, Corollary 5.1 and Conjecture 5.2.
\bibitem{BW} A. Baker and G. W\"ustholz,
\emph{Logarithmic forms and group varieties}, J. reine angew. Math. \textbf{442} (1993), 19--62.
\url{https://doi.org/10.1515/crll.1993.442.19}.
\bibitem{Nesterenko} Yu. V. Nesterenko,
\emph{Modular functions and transcendence questions}, Sb. Math. \textbf{187} (1996), 1319--1348.
\url{https://doi.org/10.1070/SM1996v187n09ABEH000158}.
\bibitem{RVY} S. Reitzes, P. Vulakh and M. P. Young,
\emph{Zeros of certain combinations of Eisenstein series},
\url{https://arxiv.org/abs/1603.01306}.
See the discussion of smaller weights $4,6,8$.
\end{thebibliography}
\end{document}
